\documentclass[superscriptaddress,amsmath, amssymb, onecolumn, pra]{revtex4-2} 

\usepackage[bookmarks=false,linkcolor=red,urlcolor=blue,colorlinks=true,citecolor=blue]{hyperref}

\pdfoutput=1 % for arXiv

\usepackage{wrapfig}
\usepackage{graphicx}  % needed for figures
\usepackage{dcolumn}   % needed for some tables
\usepackage{bm}        % for math
\usepackage{amssymb}   % for math
\usepackage{physics}

\usepackage{comment}
\usepackage{fdsymbol}
\usepackage{xcolor}
\usepackage[english]{babel}
\usepackage{amsmath,amsthm,amssymb}
\usepackage{amsfonts}
\usepackage{amsopn}
\usepackage{float}
\usepackage{url}
\usepackage{graphicx}
\usepackage{dcolumn}
\usepackage{bm}
\usepackage{color}
\hyphenation{itself}
\usepackage{float}
\usepackage{lipsum}
\usepackage[normalem]{ulem}
\usepackage[capitalize]{cleveref}
\usepackage{silence}
\WarningFilter{revtex4-2}{Repair the float}
\WarningFilter{revtex4-2}{No Society}
\usepackage{mathtools}
\DeclarePairedDelimiter\ceil{\lceil}{\rceil}
\DeclarePairedDelimiter\floor{\lfloor}{\rfloor}

\input{additional-math}

\begin{document}

\tableofcontents

\section{Fixed Temperature Spherical $p$-spin}
We consider the Spherical version of the $p$-spin model, characterized by the following Hamiltonian:
\begin{equation}\label{eq:H_pspin_spherical}
    \mathcal{H}_J(\vect{s}) = -\sum_{i_1<\dots<i_p}J_{i_1\dots i_p}s_{i_1}\dots s_{i_p}\,,\quad \vect{s}\in S_N
\end{equation}
where $J_{i_1\dots i_p}$ is a random Gaussian tensor where each entrance is i.i.d. with mean $\mathbb{E}[J_{i_1\dots i_p}]=\frac{J_0 p!}{N^{p-1}}$ and variance $\mathbb{E}[J^2_{i_1\dots i_p}] - \left(\mathbb{E}[J_{i_1\dots i_p}]\right)^2=\frac{p!}{N^{p-1}}$.
This variance makes the model the same as the one considered in \cite{alaoui_shattering_2024}.

We denote the Gibbs measure $\mu_{\beta}(\vect{s}) = \frac{1}{Z_{\beta}}\exp(-\beta\mathcal{H}_J(\vect{s})) $, with $Z_{\beta}$ the normalization constant.

We use the following definitions for $\beta_d$ and $\beta_K$:
\begin{equation}
    \beta_d(p) = \sqrt{\frac{(p-1)^{p-1}}{p (p-2)^{p-2}}}\,,\quad 
    \beta_K(p)^2 = \inf_{q \in [0,1]} \qty(\frac{1}{q^p} \log\qty(\frac{1}{1-q}) - \frac{1}{q^{p-1}}) \,,
\end{equation}
which for large values of $p$ scale as
\begin{equation}\label{eq:betad-betac-defs}
    \beta_d(p) = \sqrt{e} + O(p^{-1}) \,, \quad \beta_K(p) = \sqrt{\log p} \qty(1 + o_p(1)) \,.
\end{equation}

Our primary objective is to determine $\beta_d^+$, defined as the highest value of $\beta$ for which the system exhibits fast mixing from the worst case initial condition.
\subsection{Computation of $T_d^+$ from following states (RS)}
The RS Franz-Parisi potential can be shown to be
\begin{equation}\label{eq:FP_RS_spherical}
    -2\beta F_\beta(m,q) = 1 + \log 2\pi + \beta^2(1-q^p) + \log(1-q) + \frac{q-m^2}{1-q} + 2\beta J_0 m^p\,.
\end{equation}
where $J_0 = \beta' $, with $\beta'$ the inverse temperature of the planted solution, $m$ is the overlap with the planted solution and $q$ is the replica overlap.

\subsubsection{Finite $p$}
We first use the state following formalism from \citep{krzakala_2010,sun_2012} to get an estimate of the $T_d^+$ for the spherical case at a finite value of $p$.
We consider the planted version of the model, with the planting temperature fixed at the Kauzmann temperature $T_K$, so that we have contiguity to the random version of the model.
We start by running the fixed point equations at temperature $T = T_K$ to find the local minima with $m > 0$.
Then, we increase the temperature and follow the minimum until it vanishes, and we call $T_d^+$ the temperature value at which this happens.
For finite values of p (for $p > 10^4$, the numerical routines start to become unstable), we can look at the scaling of $T_d^+$ with p and it turns out to be consistent with the one found in the previous section, as we will show in the plots in the next section.
%In Figure \ref{fig:spherical_Tdplus_Td_Tk} we plot these values, along with the values of $T_d$ and $T_K$.

\subsubsection{Large $p$ limit}
We can rewrite the fixed-point equations of \eqref{eq:FP_RS_spherical} as following
\begin{equation}
\begin{cases}
    1-q = \frac{m^{2-p}}{\beta p J_0}\,, \\
    (1-q)^2 = \frac{(q-m^2)q^{1-p}}{p\beta^2}\,.
\end{cases}
\end{equation}
In the large $p$ limit the Kauzmann transition is at $\beta_K \sim \sqrt{\log p}$ \footnote{We use the symbol $\sim$ to indicate that two expressions are of the same order in $p$.}\cite{kurchan1993barriers} and, following the numerical results, we consider $\beta = \frac{C_+}{\sqrt{\log p}}$, with $C_+$ is a constant to be determined.

Then from the first fixed point we get $1-q \sim (pm^p)^{-1}$. Numerically, one can verify that $m^p$ converges for $p$ large to a positive constant, such that we can say $q = 1 - \frac{\alpha_q}{p} + o(\frac{1}{p})$, and thus $q^p\to e^{-\alpha_q}$.

Putting this scaling in the second fixed point, we get
\begin{equation}
    \frac{\alpha_q}{p} \sim \sqrt{\frac{\log p}{p}(q-m^2)}\,,
\end{equation}
which in turn implies $q-m^2 = \frac{\gamma}{p\log p}$. By a simple argument, one can see that this means $m = 1 - \frac{\alpha_m}{p}= 1 - \frac{\alpha_q}{2p} + o(\frac{1}{p})$ and thus that $m^p\to e^{-\alpha_q/2}$.

Directly by imposing these scaling back into \eqref{eq:FP_RS_spherical} one finds the following relationships between the constants:
\begin{equation}
    \begin{cases}
        \frac{C_+}{2}\alpha_q = e^{\frac{\alpha_q}{2}}\,, \\
        C_+^2\alpha_q^2 = 4\gamma e^{\alpha_q}\,.
    \end{cases}
\end{equation}
Using the first equation to specify $C_+$ in terms of $\alpha_q$ and substituting in the second equation we find $\gamma = 1/2$ and we are left with one equation in two variables.

We derive the last equation by looking at the Hessian of the free energy $F_{\beta}(m,q)$. Indeed we expect that $T_d^+$ is the temperature at which the local minima disappears, and thus it is also where the smaller eigenvalue of the Hessian becomes zero. 
The eigenvalues can be computed and the smaller one is given by

\begin{equation}
    \begin{cases}
        \lambda_{\rm min} = -\frac{1}{4 \beta  m^2 (q-1)^3 q^2}\left(A - \sqrt{A^2 + B}\right)\,, \\
        \begin{aligned}
        A = 2 \beta  J_0 (p-1) p (q-1)^3 q^2 m^p+2 m^4 q^2 -m^2 \left(\beta ^2 (p-1) p (q-1)^3 q^p- q^2 (q (2 q-5)+1)\right) \,, 
        \end{aligned}\\
        \begin{aligned}
        B = 8&m^2 (q-1)^2 q^2 \Bigg(\beta  J_0 (p-1) p (q-1) m^p \left( q^2 \left(-2 m^2+q+1\right)+\beta ^2 (p-1) p (q-1)^3
   q^p\right)\\ &+m^2 \left(\beta ^2 (p-1) p (q-1)^3 q^p+ (q+1) q^2\right)\Bigg)\,.
        \end{aligned}
    \end{cases}
\end{equation}
One can check that $B = o(A)$ and thus it is easy to convince ourselves that $B = 0 \Rightarrow\lambda_{\min} = 0$. At the first order in $B$, this is equivalent to asking 
\begin{equation}
    m^2q^2(1+q) + m^pp^2(q-1)q^2\beta J_0(1-2m^2 + q) = 0 \,.
\end{equation}
After some trivial computations, one arrives to the following formula:
\begin{equation}
    C_+\alpha_q^2 = 2e^{\frac{\alpha_q}{2}} \,,
\end{equation}
which combined with the one coming from the fixed point equations leads to $\alpha_q = 2$, and thus $\alpha_m=1$ and $C_+ = \frac{e}{2}$.

We verify these scalings numerically; in Fig. \ref{fig:spherical2_Tdplus_Td_Tk_log} for $T_d^+$ and in Fig. \ref{fig:spherical2_m_q} for $m$ and $q$.
\begin{figure}
    \centering
    \includegraphics[width=0.75\linewidth]{Figures/spherical2_Tdplus_TkTd_log.png}
    \caption{Numerical estimation of $T_d^+$ from state following, superimposed with the asymptotic scalings.}
    \label{fig:spherical2_Tdplus_Td_Tk_log}
\end{figure}
\begin{figure}
    \centering
    \includegraphics[width=1\linewidth]{Figures/spherical2_m_q.png}
    \caption{Check of the asymptotic scalings of $m$ and $q$.}
    \label{fig:spherical2_m_q}
\end{figure}
\section{Fixed Temperature Ising $p$-spin}
We come back to the Ising version of the $p$-spin model, characterized by the Hamiltonian 
\begin{equation}\label{eq:H_pspin_spherical}
    \mathcal{H}_J(\vect{s}) = -\sum_{i_1<\dots<i_p}J_{i_1\dots i_p}s_{i_1}\dots s_{i_p}\,,\quad s_i\in \{-1,1\}\;\forall\,i\in[1:N]
\end{equation}
where $J_{i_1\dots i_p}$ is a random Gaussian tensor where each entrance is i.i.d. with mean $\mathbb{E}[J_{i_1\dots i_p}]=\frac{J_0 p!}{N^{p-1}}$ and variance $\mathbb{E}[J^2_{i_1\dots i_p}] - \left(\mathbb{E}[J_{i_1\dots i_p}]\right)^2=\frac{p!}{N^{p-1}}$.
This variance makes the model the same as the one considered in \cite{alaoui_shattering_2024}.

We denote the Gibbs measure $\mu_{\beta}(\vect{s}) = \frac{1}{Z_{\beta}}\exp(-\beta\mathcal{H}_J(\vect{s})) $, with $Z_{\beta}$ the normalization constant.
The RS Franz-Parisi potential can be computed through the cavity method \cite{krzakala_2010} and turns out to be {\color{red} To Add}.
\subsection{Computation of $T_d^+$ from following states}
\subsubsection{Finite $p$}
In this section, we use again the state following formalism from {\color{red} Add ref} to get an estimate of the $T_d^+$ for the Ising case, in the same we did for the spherical version.
\begin{figure}
    \centering
    \includegraphics[width=0.75\linewidth]{Figures/ising_Tdplus_TkTd_log.png}
    \caption{Numerical estimation of $T_d^+$ from state following for the Ising p-spin model.}
    \label{fig:ising_Tdplus_Td_Tk}
\end{figure}
In Figure \ref{fig:ising_Tdplus_Td_Tk} we plot the estimated values of $T_d^+$, along with the values of $T_d$ and $T_K$.

\appendix


 \bibliographystyle{alpha} % We choose the "plain" reference style
 \bibliography{fixed-energy}

\end{document}